% Compile: xelatex main.tex (รัน 2 รอบ)
\documentclass[12pt,a4paper]{article}

\usepackage{fontspec}
\usepackage{polyglossia}
\setdefaultlanguage{thai}
\setotherlanguage{english}

\setmainfont{Sarabun}[
    BoldFont       = Sarabun-Bold,
    ItalicFont     = Sarabun-Italic,
    BoldItalicFont = Sarabun-BoldItalic,
]
\setmonofont{Consolas}[Scale=0.92]
\newfontfamily\thaifonttt{Consolas}[Scale=0.92]

\usepackage{amsmath, amssymb}
\usepackage{graphicx}
\usepackage{booktabs}
\usepackage{geometry}
\usepackage{xcolor}
\usepackage{float}
\usepackage{enumitem}
\usepackage{tcolorbox}
\usepackage{needspace}
\usepackage{caption}
\usepackage{listings}
\usepackage{upquote}
\tcbuselibrary{skins, breakable}

\XeTeXlinebreaklocale "th"
\XeTeXlinebreakskip = 0pt plus 1pt
\sloppy

\definecolor{captioncol}{RGB}{80, 80, 80}
\definecolor{accentblue}{RGB}{37, 99, 235}
\definecolor{codebg}{RGB}{248, 250, 252}
\definecolor{codeframe}{RGB}{203, 213, 225}
\definecolor{keyword}{RGB}{124, 58, 237}
\definecolor{string}{RGB}{22, 163, 74}
\definecolor{comment}{RGB}{100, 116, 139}
\definecolor{output}{RGB}{30, 58, 138}
\captionsetup{font={small, color=captioncol}, justification=centering}

\lstdefinestyle{pycode}{
    language=Python,
    backgroundcolor=\color{codebg},
    basicstyle=\ttfamily\small,
    keywordstyle=\color{keyword}\bfseries,
    stringstyle=\color{string},
    commentstyle=\color{comment}\itshape,
    numberstyle=\tiny\color{comment},
    numbers=left,
    numbersep=8pt,
    frame=single,
    rulecolor=\color{codeframe},
    breaklines=true,
    breakatwhitespace=false,
    showstringspaces=false,
    tabsize=4,
    xleftmargin=16pt,
    xrightmargin=4pt,
    columns=flexible,
    keepspaces=true,
    literate={≈}{{$\approx$}}1 {α}{{$\alpha$}}1 {β}{{$\beta$}}1
}

\lstdefinestyle{output}{
    basicstyle=\ttfamily\small\color{output},
    backgroundcolor=\color{codebg},
    frame=single,
    rulecolor=\color{codeframe},
    numbers=none,
    xleftmargin=8pt,
    xrightmargin=4pt,
    breaklines=true,
    breakatwhitespace=false,
    columns=flexible,
    keepspaces=true,
    literate={≈}{{$\approx$}}1 {α}{{$\alpha$}}1 {β}{{$\beta$}}1
}

\widowpenalty=10000
\clubpenalty=10000

\geometry{top=2.5cm, bottom=2.5cm, left=2.5cm, right=2.5cm, headheight=14pt}

\newtcolorbox{ansbox}[1][]{
    colback=green!5!white, colframe=green!50!black,
    boxrule=0.5pt, arc=4pt, left=8pt, right=8pt, top=5pt, bottom=5pt, #1
}

% =================================================================
\begin{document}

\begin{center}
    {\Large\textbf{CPE 342 Machine Learning}}\\[2pt]
    {\Large\textbf{Assignment 1: Training Models}}\\[4pt]
    {\large OLS Regression}\\[6pt]
    \textcolor{accentblue}{67070501042 วิศิษฐ์ สุวรรณเนาว์}
\end{center}

\noindent\rule{\textwidth}{0.4pt}\vspace{4pt}

\noindent
ข้อมูลที่กำหนดให้เป็นตัวเลขงบประมาณโฆษณา ($X$, หน่วย: พันดอลลาร์) และยอดขาย ($Y$, หน่วย: พันหน่วย) จาก 10 เดือน ดังนี้

\begin{table}[H]
\centering
\begin{tabular}{c c c}
\toprule
\textbf{Month} & \textbf{Advertising Budget $X$ (\$1,000s)} & \textbf{Sales $Y$ (1,000 units)} \\
\midrule
1  & 3.0  & 6.0  \\
2  & 5.0  & 9.0  \\
3  & 2.0  & 4.0  \\
4  & 7.0  & 10.0 \\
5  & 8.0  & 12.0 \\
6  & 1.0  & 3.0  \\
7  & 4.0  & 7.0  \\
8  & 6.0  & 8.0  \\
9  & 9.0  & 13.0 \\
10 & 10.0 & 15.0 \\
\bottomrule
\end{tabular}
\end{table}

% =================================================================
\section*{Task 1: Implement OLS Regression}
% =================================================================

\subsection*{1a) Fit Simple Linear Regression ด้วยวิธี OLS}

\textbf{แบบจำลอง:}
สมมติว่าความสัมพันธ์ระหว่าง $X$ กับ $Y$ เป็นเชิงเส้น (linear) กำหนดให้ $\hat{Y}_i = \alpha + \beta X_i$ 
โดยที่ $\alpha$ คือ intercept และ $\beta$ คือ slope 

แทน $Y_i = \beta X_i + \alpha$ ลงในสมการ error $\varepsilon_i = Y_i - \hat{Y}_i$ จะได้
\[
    \varepsilon_i = Y_i - (\alpha + \beta X_i)
\]

\medskip\noindent
\textbf{หลักการของ OLS:}
ฟังก์ชันผลรวมของกำลังสองของ error (Least Squares Function) $S$:
\[
    S(\alpha, \beta) = \sum_{i=1}^{n} \varepsilon_i^2 = \sum_{i=1}^{n} \bigl(Y_i - \alpha - \beta X_i\bigr)^2
\]
ฟังก์ชัน $S$ จะมีค่าน้อยที่สุดเมื่อ $\frac{\partial S}{\partial \alpha} = 0$ และ $\frac{\partial S}{\partial \beta} = 0$

\medskip\noindent
\textbf{การหา partial derivatives:}

\begin{align}
    \frac{\partial S}{\partial \alpha} &= -2\sum_{i=1}^{n}(Y_i - \alpha - \beta X_i) = 0 \label{eq:da}\\[6pt]
    \frac{\partial S}{\partial \beta}  &= -2\sum_{i=1}^{n}(Y_i - \alpha - \beta X_i)X_i = 0 \label{eq:db}
\end{align}

จากสมการ \eqref{eq:da}: 
$\displaystyle\sum_{i=1}^{n} Y_i - n\alpha - \beta\sum_{i=1}^{n} X_i = 0 \implies$ \textbf{Normal Equation 1:}
\[
    n\alpha + \beta\sum_{i=1}^{n} X_i = \sum_{i=1}^{n} Y_i
\]

จากสมการ \eqref{eq:db}: 
$\displaystyle\sum_{i=1}^{n} X_i Y_i - \alpha\sum_{i=1}^{n} X_i - \beta\sum_{i=1}^{n} X_i^2 = 0 \implies$ \textbf{Normal Equation 2:}
\[
    \alpha\sum_{i=1}^{n} X_i + \beta\sum_{i=1}^{n} X_i^2 = \sum_{i=1}^{n} X_i Y_i
\]

สามารถเขียน Normal Equation ในรูปแบบ Matrix ได้ดังนี้:
\[
    \begin{bmatrix}
        n & \sum X_i \\
        \sum X_i & \sum X_i^2
    \end{bmatrix}
    \begin{bmatrix}
        \alpha \\
        \beta
    \end{bmatrix}
    =
    \begin{bmatrix}
        \sum Y_i \\
        \sum X_i Y_i
    \end{bmatrix}
\]

\medskip\noindent
\textbf{คำนวณค่า summation จากข้อมูล:}

\begin{table}[H]
\centering
\begin{tabular}{c r r r r}
\toprule
$i$ & $X_i$ & $Y_i$ & $X_i Y_i$ & $X_i^2$ \\
\midrule
1  & 3.0  & 6.0  & 18.0   & 9.0   \\
2  & 5.0  & 9.0  & 45.0   & 25.0  \\
3  & 2.0  & 4.0  & 8.0    & 4.0   \\
4  & 7.0  & 10.0 & 70.0   & 49.0  \\
5  & 8.0  & 12.0 & 96.0   & 64.0  \\
6  & 1.0  & 3.0  & 3.0    & 1.0   \\
7  & 4.0  & 7.0  & 28.0   & 16.0  \\
8  & 6.0  & 8.0  & 48.0   & 36.0  \\
9  & 9.0  & 13.0 & 117.0  & 81.0  \\
10 & 10.0 & 15.0 & 150.0  & 100.0 \\
\midrule
$\sum$ & \textbf{55.0} & \textbf{87.0} & \textbf{583.0} & \textbf{385.0} \\
\bottomrule
\end{tabular}
\end{table}

ดังนั้น:
\[
    n = 10, \quad
    \sum X_i = 55, \quad
    \sum Y_i = 87, \quad
    \sum X_i Y_i = 583, \quad
    \sum X_i^2 = 385
\]

\begin{ansbox}
\textbf{สรุป 1a:} \quad จากข้อมูลข้างต้น จะได้ระบบสมการ:
\[
    \begin{bmatrix}
        10 & 55 \\
        55 & 385
    \end{bmatrix}
    \begin{bmatrix}
        \alpha \\
        \beta
    \end{bmatrix}
    =
    \begin{bmatrix}
        87 \\
        583
    \end{bmatrix}
\]
\end{ansbox}

% ----- 1b -----
\subsection*{1b) คำนวณ Regression Coefficients}

เราสามารถหาค่า $\alpha$ (intercept) และ $\beta$ (slope) ได้หลายวิธี:

\subsubsection*{วิธีที่ 0: สูตร $S_{xx}$, $S_{xy}$ (Closed-form จาก Normal Equations)}
จาก Normal Equations สามารถแก้หาค่า $\beta$ ได้โดยตรงด้วยการนิยามทั้งสองแบบ:
\begin{align*}
    S_{xx} &= \sum_{i=1}^{n}(X_i - \bar{X})^2 = \sum X_i^2 - \frac{(\sum X_i)^2}{n} = 385 - \frac{55^2}{10} = 82.5 \\[4pt]
    S_{xy} &= \sum_{i=1}^{n}(X_i - \bar{X})(Y_i - \bar{Y}) = \sum X_iY_i - \frac{\sum X_i \sum Y_i}{n} = 583 - \frac{(55)(87)}{10} = 104.5
\end{align*}
จะได้:
\[
    \beta = \frac{S_{xy}}{S_{xx}} = \frac{104.5}{82.5} \approx 1.2667, \qquad \alpha = \bar{Y} - \beta\bar{X} = 8.7 - (1.2667)(5.5) \approx 1.7333
\]
วิธีนี้ให้ผลลัพธ์ที่ตรงกันกับทุกวิธีอื่น แต่คำนวณได้เร็วกว่าเพราะไม่ต้องแก้ระบบสมการ

\subsubsection*{วิธีที่ 1: Cramer's Rule}
จากระบบสมการ:
\begin{align*}
    10\alpha + 55\beta &= 87 \\
    55\alpha + 385\beta &= 583
\end{align*}
หาค่า Determinant:
\[
    \Delta = \det \begin{bmatrix} 10 & 55 \\ 55 & 385 \end{bmatrix} = (10)(385) - (55)(55) = 3850 - 3025 = 825
\]
\[
    \Delta_\alpha = \det \begin{bmatrix} 87 & 55 \\ 583 & 385 \end{bmatrix} = (87)(385) - (55)(583) = 33495 - 32065 = 1430
\]
\[
    \Delta_\beta = \det \begin{bmatrix} 10 & 87 \\ 55 & 583 \end{bmatrix} = (10)(583) - (87)(55) = 5830 - 4785 = 1045
\]
จะได้:
\[
    \alpha = \frac{\Delta_\alpha}{\Delta} = \frac{1430}{825} = \frac{26}{15} \approx 1.7333
\]
\[
    \beta = \frac{\Delta_\beta}{\Delta} = \frac{1045}{825} = \frac{19}{15} \approx 1.2667
\]

\subsubsection*{วิธีที่ 2: Pseudoinverse Method}
กำหนด $A = \begin{bmatrix} 1 & X_1 \\ 1 & X_2 \\ \vdots & \vdots \\ 1 & X_n \end{bmatrix}$, 
$Y = \begin{bmatrix} Y_1 \\ Y_2 \\ \vdots \\ Y_n \end{bmatrix}$, 
และ $c = \begin{bmatrix} \alpha \\ \beta \end{bmatrix}$

จากสมการ $Ac = Y$ เราใช้ Left Pseudoinverse:
\[
    c = (A^T A)^{-1} A^T Y
\]
จากตารางคำนวณ:
\[
    A^T A = \begin{bmatrix} n & \sum X_i \\ \sum X_i & \sum X_i^2 \end{bmatrix} = \begin{bmatrix} 10 & 55 \\ 55 & 385 \end{bmatrix}
\]
\[
    A^T Y = \begin{bmatrix} \sum Y_i \\ \sum X_i Y_i \end{bmatrix} = \begin{bmatrix} 87 \\ 583 \end{bmatrix}
\]
หา $(A^T A)^{-1}$:
\[
    (A^T A)^{-1} = \frac{1}{825} \begin{bmatrix} 385 & -55 \\ -55 & 10 \end{bmatrix}
\]
คำนวณหา $c$:
\begin{align*}
    c &= \frac{1}{825} \begin{bmatrix} 385 & -55 \\ -55 & 10 \end{bmatrix} \begin{bmatrix} 87 \\ 583 \end{bmatrix} \\[6pt]
      &= \frac{1}{825} \begin{bmatrix} (385)(87) + (-55)(583) \\ (-55)(87) + (10)(583) \end{bmatrix} \\[6pt]
      &= \frac{1}{825} \begin{bmatrix} 33495 - 32065 \\ -4785 + 5830 \end{bmatrix} \\[6pt]
      &= \frac{1}{825} \begin{bmatrix} 1430 \\ 1045 \end{bmatrix} = \begin{bmatrix} 1430/825 \\ 1045/825 \end{bmatrix} = \begin{bmatrix} 26/15 \\ 19/15 \end{bmatrix}
\end{align*}

\begin{ansbox}
\textbf{สรุป 1b:} \quad ได้ค่าสัมประสิทธิ์ดังนี้:
\begin{itemize}[noitemsep, topsep=2pt]
    \item \textbf{Intercept ($\alpha$)} = $\frac{26}{15} \approx 1.7333$
    \item \textbf{Slope ($\beta$)} = $\frac{19}{15} \approx 1.2667$
\end{itemize}
\end{ansbox}
\newpage

% ----- 1c -----
\subsection*{1c) สมการ Fitted Line}

นำค่า $\alpha$ และ $\beta$ ที่คำนวณได้ไปแทนในสมการ $\hat{Y} = \alpha + \beta X$

\begin{ansbox}
\textbf{สรุป 1c:} \quad สมการเส้นตรงพยากรณ์คือ \quad $\boxed{\hat{Y} = 1.7333 + 1.2667X}$
\end{ansbox}

\begin{figure}[H]
    \centering
    \includegraphics[width=0.75\textwidth]{plot_1_regression.pdf}
    \caption{OLS Linear Regression Analysis (Advertising Budget vs Sales) แสดงข้อมูลจริงและเส้นพยากรณ์ $\hat{Y} = 1.7333 + 1.2667X$}
\end{figure}
\newpage

% =================================================================
\section*{Task 2: Interpret the Results}
% =================================================================

\subsection*{2a) ความหมายของ Slope และ Intercept}

\textbf{ความหมายของ Slope ($\beta \approx 1.2667$):}
\begin{ansbox}
Slope บ่งบอกว่ายอดขาย (Sales) ถูกคาดการณ์ว่าจะเพิ่มตามค่าโฆษณา (Advertising) ที่เพิ่มขึ้น นั่นคือ ทุก ๆ \$1,000 ของ advertising ที่เพิ่มขึ้น จะทำให้ยอดขาย (Sales) เพิ่มขึ้นประมาณ \textbf{1,267 หน่วย} ค่านี้เป็นความสัมพันธ์เชิงสถิติ (association) ไม่ใช่การพิสูจน์ว่าโฆษณา\textit{เป็นสาเหตุ}(causation) ของยอดขายที่เพิ่มขึ้น
\end{ansbox}

\textbf{ความหมายของ Intercept ($\alpha \approx 1.7333$):}
\begin{ansbox}
จุดตัดแกน Y บ่งบอกยอดขายที่คาดการณ์ไว้เมื่อไม่มีต้นทุนค่าโฆษณา กล่าวคือ ถ้าบริษัทไม่ทำโฆษณาเลย ($X=0$) จะสามารถคาดการณ์ยอดขายได้ประมาณ \textbf{1,733 หน่วย}
\end{ansbox}


\subsection*{2b) Predict ยอดขายเมื่องบโฆษณา = \$12,000}

แทน $X = 12$ ลงในสมการที่ได้:
\begin{align*}
    \hat{Y} &= 1.7333 + 1.2667(12) \\
    \hat{Y} &= 1.7333 + 15.2004 \\
    \hat{Y} &\approx 16.9333
\end{align*}

\begin{ansbox}
\textbf{สรุป 2b:} \quad ถ้าต้นทุนในการโฆษณาคือ \$12,000 คาดการณ์ยอดขายได้ประมาณ \textbf{16.93 พันหน่วย} (ประมาณ 16,933 หน่วย)

\textit{ข้อควรระวัง:} $X = 12$ อยู่นอกช่วงข้อมูลที่ใช้ fit โมเดล (ค่าสูงสุดใน dataset คือ $X = 10$) การทำนายนี้จึงเป็น \textbf{extrapolation} ซึ่งโมเดลไม่สามารถยืนยันได้ว่าความสัมพันธ์เชิงเส้นยังคงดำเนินต่อไปที่ระดับนี้
\end{ansbox}

\begin{figure}[H]
    \centering
    \includegraphics[width=0.65\textwidth]{plot_7_actual_vs_pred.pdf}
    \caption{Actual vs Predicted Sales แสดงความแม่นยำของการพยากรณ์โดยเปรียบเทียบค่าพยากรณ์กับค่าจริง}
\end{figure}
\newpage

% =================================================================
\section*{Task 3: Model Evaluation}
% =================================================================

\subsection*{3a) คำนวณ R-squared ($R^2$)}

$R^2$ วัดความเหมาะสมของแบบจำลอง (Goodness of fit) นิยามเป็น:
\[
    R^2 = 1 - \frac{SS_{\text{res}}}{SS_{\text{tot}}}
\]
โดย:
\begin{itemize}[noitemsep, topsep=4pt]
    \item $SS_{\text{res}} = \displaystyle\sum(Y_i - \hat{Y}_i)^2$ 
    \item $SS_{\text{tot}} = \displaystyle\sum(Y_i - \bar{Y})^2$ 
\end{itemize}

\textbf{คำนวณ:} (โดย $\bar{Y} = 8.7$)

\begin{table}[H]
\centering
\small
\begin{tabular}{c r r r r r r}
\toprule
$i$ & $X_i$ & $Y_i$ & $\hat{Y}_i = 1.7333 + 1.2667X_i$ & $e_i = Y_i - \hat{Y}_i$ & $e_i^2$ & $(Y_i - \bar{Y})^2$ \\
\midrule
1  & 3.0  & 6.0  & 5.5333  & $0.4667$  & 0.2178  & 7.2900   \\
2  & 5.0  & 9.0  & 8.0667  & $0.9333$  & 0.8711  & 0.0900   \\
3  & 2.0  & 4.0  & 4.2667  & $-0.2667$ & 0.0711  & 22.0900  \\
4  & 7.0  & 10.0 & 10.6000 & $-0.6000$ & 0.3600  & 1.6900   \\
5  & 8.0  & 12.0 & 11.8667 & $0.1333$  & 0.0178  & 10.8900 \\
6  & 1.0  & 3.0  & 3.0000  & $0.0000$  & 0.0000  & 32.4900  \\
7  & 4.0  & 7.0  & 6.8000  & $0.2000$  & 0.0400  & 2.8900   \\
8  & 6.0  & 8.0  & 9.3333  & $-1.3333$ & 1.7778  & 0.4900   \\
9  & 9.0  & 13.0 & 13.1333 & $-0.1333$ & 0.0178  & 18.4900 \\
10 & 10.0 & 15.0 & 14.4000 & $0.6000$  & 0.3600  & 39.6900 \\
\midrule
$\sum$ & & & & & $\mathbf{3.7333}$ & $\mathbf{136.1000}$ \\
\bottomrule
\end{tabular}
\end{table}

จากตาราง: $SS_{\text{res}} = 3.7333$, $SS_{\text{tot}} = 136.1000$

แทนค่า:
\[
    R^2 = 1 - \frac{3.7333}{136.10} = 1 - 0.0274 = 0.9726
\]

\begin{ansbox}
\textbf{สรุป 3a:} \quad $\boxed{R^2 \approx 0.9726}$ หรือประมาณ 97.26\%
\end{ansbox}

\subsection*{3b) ตีความค่า R-squared}

\begin{ansbox}
\textbf{สรุป 3b:} \quad ค่า R-squared ที่ $0.9726$ หมายความว่า ประมาณ \textbf{97.26\% ของการเปลี่ยนแปลงของยอดขาย (Sales) สามารถถูกอธิบายได้ด้วยงบโฆษณา (Advertising Budget)}

จากตาราง residuals ข้างต้น พบว่า $\bar{e} = 0$ \textit{(residual mean เท่ากับศูนย์พอดี แสดงว่าโมเดลไม่มี bias)} โดย Standard Deviation ของ residuals อยู่ที่ $s_e = 0.6441$ พันหน่วย (เทียบกับค่าเฉลี่ยยได้ดี เพราะ residual สูงสุดใน dataset คือ $|e_8| = 1.33$)

\textit{สิ่งที่ $R^2$ สูงไม่ได้บ่งชี้:} ว่างบโฆษณา\emph{เป็นสาเหตุ}ของยอดขาย หรือรับประกันว่าโมเดลจะพยากรณ์ได้ดีเสมอในอนาคต
\end{ansbox}

\subsection*{Diagnostic Plots (วิเคราะห์ความสัมพันธ์เชิงลึก)}

\begin{figure}[H]
    \centering
    \includegraphics[width=0.75\textwidth]{plot_2_residuals_lollipop.pdf}
    \caption{\textbf{Residuals Lollipop Chart:} แสดงขนาดของความคลาดเคลื่อน (Residual) ของแต่ละจุดข้อมูล (M1-M10) ช่วยให้เห็นชัดเจนว่าจุดไหนที่โมเดลทำนายคลาดเคลื่อนมากที่สุด (เช่น M8 มีค่า Residual ต่ำกว่าศูนย์มากที่สุด)}
\end{figure}

\begin{figure}[H]
    \centering
    \includegraphics[width=0.75\textwidth]{plot_3_res_vs_fitted.pdf}
    \caption{\textbf{Residuals vs Fitted Values:} ใช้ตรวจสอบรูปแบบความคลาดเคลื่อน หากจุดกระจายตัวแบบสุ่มรอบๆ เส้น 0 แสดงว่าโมเดล Linear เหมาะสม ไม่มีปัญหา non-linearity}
\end{figure}

\begin{figure}[H]
    \centering
    \includegraphics[width=0.75\textwidth]{plot_4_res_dist.pdf}
    \caption{\textbf{Residuals Distribution:} ตรวจสอบการกระจายตัวของความคลาดเคลื่อนว่าใกล้เคียงการแจกแจงแบบปกติ (Normal Distribution) หรือไม่ ซึ่งเป็นสมมติฐานสำคัญของ OLS}
\end{figure}

\begin{figure}[H]
    \centering
    \includegraphics[width=0.75\textwidth]{plot_5_qq.pdf}
    \caption{\textbf{Normal Q-Q Plot:} ใช้เปรียบเทียบการกระจายตัวของ Residuals กับ Normal Distribution ในเชิง Quantile หากจุดส่วนใหญ่อยู่บนเส้นทแยงมุม แสดงว่า Residuals มีการแจกแจงแบบปกติ}
\end{figure}

\begin{figure}[H]
    \centering
    \includegraphics[width=0.75\textwidth]{plot_6_scale_loc.pdf}
    \caption{\textbf{Scale-Location Plot:} ใช้ตรวจสอบสมมติฐาน Homoscedasticity (ความแปรปรวนคงที่) หากจุดกระจายตัวสม่ำเสมอโดยไม่มีรูปแบบที่ชัดเจน แสดงว่าความแปรปรวนของความคลาดเคลื่อนคงที่ตลอดทุกช่วงค่าพยากรณ์}
\end{figure}
\clearpage

% =================================================================
\section*{สรุปผลลัพธ์}
% =================================================================

\begin{table}[H]
\centering
\begin{tabular}{l l}
\toprule
\textbf{ข้อ} & \textbf{คำตอบ} \\
\midrule
\textbf{1a} & ระบบสมการ $10\alpha + 55\beta = 87$ และ $55\alpha + 385\beta = 583$ \\
\textbf{1b} & Intercept ($\alpha$) $\approx 1.7333$, \quad Slope ($\beta$) $\approx 1.2667$ \\
\textbf{1c} & สมการ Fitted Line: $\hat{Y} = 1.7333 + 1.2667X$ \\
\textbf{2a} & Slope: เพิ่มงบโฆษณา \$1,000 ทำให้ยอดขายเพิ่ม 1,267 หน่วย \\
            & Intercept: หากไม่มีงบโฆษณาเลย ยอดขายคาดการณ์คือ 1,733 หน่วย \\
\textbf{2b} & เมื่องบโฆษณา = \$12,000 คาดการณ์ยอดขายได้ประมาณ 16,933 หน่วย \\
\textbf{3a} & $R^2 \approx 0.9726$ (97.26\%) \\
\textbf{3b} & 97.26\% ของยอดขายสามารถอธิบายได้ด้วยงบโฆษณา (โมเดลมีความเหมาะสมดีเยี่ยม) \\
\bottomrule
\end{tabular}
\end{table}

\newpage
% =================================================================
\section*{Appendix: Jupyter Notebook Output}
% =================================================================

\noindent
ส่วนนี้คือโค้ด Python แบบเต็มและผลลัพธ์ (Output) ที่รันจริงจาก \texttt{Assignment\_1\_OLS.ipynb} เพื่อแสดงวิธีการคำนวณและสร้างกราฟทั้งหมด

% Auto-generated notebook appendix
\begin{tcolorbox}[colback=blue!5!white,colframe=blue!75!black,title=\textbf{Jupyter Notebook: Assignment\_1\_OLS.ipynb}]
โค้ดและผลลัพธ์ทั้งหมดด้านล่างเป็นการรันจริงจาก Jupyter Notebook
\end{tcolorbox}

\vspace{1em}
\begin{tcolorbox}[colback=gray!5!white,colframe=gray!50!black,boxrule=0.5pt,arc=2pt,left=2pt,right=2pt,top=2pt,bottom=2pt]
\textbf{\large 0. Imports and Setup}\par\smallskip
\end{tcolorbox}
\vspace{1em}
\needspace{5\baselineskip}
\noindent\textbf{In [1]:}
\begin{lstlisting}[style=pycode]
%matplotlib inline
import matplotlib as mpl
mpl.rcParams['figure.dpi'] = 72
import numpy as np
import matplotlib
import matplotlib.pyplot as plt
import matplotlib.font_manager as fm
import os
import scipy.stats as stats

# --- Robust Sarabun Font Setup ---
font_dir = os.path.join(os.environ.get('LOCALAPPDATA', ''), 'Microsoft', 'Windows', 'Fonts')
sarabun_r_path = os.path.join(font_dir, 'Sarabun-Regular.ttf')
sarabun_b_path = os.path.join(font_dir, 'Sarabun-Bold.ttf')

prop_r = fm.FontProperties(fname=sarabun_r_path, size=11)
prop_b = fm.FontProperties(fname=sarabun_b_path, size=12)
prop_title = fm.FontProperties(fname=sarabun_b_path, size=14)
prop_small = fm.FontProperties(fname=sarabun_r_path, size=9)

def apply_font(ax):
    for label in (ax.get_xticklabels() + ax.get_yticklabels()):
        label.set_fontproperties(prop_r)

plt.style.use('seaborn-v0_8-whitegrid')
print("Plotting libraries and fonts loaded.")
\end{lstlisting}
\smallskip
\needspace{5\baselineskip}
\noindent\textbf{Out[1]:}
\begin{lstlisting}[style=output]
Plotting libraries and fonts loaded.
\end{lstlisting}
\vspace{1em}
\begin{tcolorbox}[colback=gray!5!white,colframe=gray!50!black,boxrule=0.5pt,arc=2pt,left=2pt,right=2pt,top=2pt,bottom=2pt]
\textbf{\large 1. Data}\par\smallskip
\end{tcolorbox}
\vspace{1em}
\needspace{5\baselineskip}
\noindent\textbf{In [2]:}
\begin{lstlisting}[style=pycode]
# Advertising budget (X, in $1,000s) and Sales (Y, in 1,000 units)
X = np.array([3., 5., 2., 7., 8., 1., 4., 6., 9., 10.])
Y = np.array([6., 9., 4., 10., 12., 3., 7., 8., 13., 15.])
n = len(X)
months = np.arange(1, n + 1)

print(f"n = {n}")
print(f"X = {X}")
print(f"Y = {Y}")
\end{lstlisting}
\smallskip
\needspace{5\baselineskip}
\noindent\textbf{Out[2]:}
\begin{lstlisting}[style=output]
n = 10
X = [ 3.  5.  2.  7.  8.  1.  4.  6.  9. 10.]
Y = [ 6.  9.  4. 10. 12.  3.  7.  8. 13. 15.]
\end{lstlisting}
\vspace{1em}
\begin{tcolorbox}[colback=gray!5!white,colframe=gray!50!black,boxrule=0.5pt,arc=2pt,left=2pt,right=2pt,top=2pt,bottom=2pt]
\textbf{\large 2. Task 1: OLS Regression}\par\smallskip
\textbf{2a. Setting up the Normal Equations}\par\smallskip
Minimise $S(\alpha, \beta) = \sum_i (Y_i - \alpha - \beta X_i)^2$ gives the Normal Equations:
$$\begin{bmatrix} n & \sum X_i \\ \sum X_i & \sum X_i^2 \end{bmatrix} \begin{bmatrix} \alpha \\ \beta \end{bmatrix} = \begin{bmatrix} \sum Y_i \\ \sum X_i Y_i \end{bmatrix}$$
\end{tcolorbox}
\vspace{1em}
\needspace{5\baselineskip}
\noindent\textbf{In [3]:}
\begin{lstlisting}[style=pycode]
sum_X   = X.sum()
sum_Y   = Y.sum()
sum_XY  = (X * Y).sum()
sum_X2  = (X ** 2).sum()

print("Summation table:")
print(f"  sum(X)  = {sum_X}")
print(f"  sum(Y)  = {sum_Y}")
print(f"  sum(XY) = {sum_XY}")
print(f"  sum(X2) = {sum_X2}")
print()
print("Normal Equation (matrix form):")
print(f"  [{n:2d}  {sum_X:5.1f}] [alpha]   [{sum_Y:5.1f}]")
print(f"  [{sum_X:2.0f} {sum_X2:5.1f}] [beta ] = [{sum_XY:5.1f}]")
\end{lstlisting}
\smallskip
\needspace{5\baselineskip}
\noindent\textbf{Out[3]:}
\begin{lstlisting}[style=output]
Summation table:
  sum(X)  = 55.0
  sum(Y)  = 87.0
  sum(XY) = 583.0
  sum(X2) = 385.0

Normal Equation (matrix form):
  [10   55.0] [alpha]   [ 87.0]
  [55 385.0] [beta ] = [583.0]
\end{lstlisting}
\vspace{1em}
\begin{tcolorbox}[colback=gray!5!white,colframe=gray!50!black,boxrule=0.5pt,arc=2pt,left=2pt,right=2pt,top=2pt,bottom=2pt]
\textbf{2b. Solving for Coefficients}\par\smallskip
\textbf{Method 0 — Closed-form via $S_{xx}$, $S_{xy}$} (derived directly from Normal Equations):
$$\beta = \frac{S_{xy}}{S_{xx}}, \qquad \alpha = \bar{Y} - \beta \bar{X}$$
\end{tcolorbox}
\vspace{1em}
\needspace{5\baselineskip}
\noindent\textbf{In [4]:}
\begin{lstlisting}[style=pycode]
X_bar = X.mean()   # 5.5
Y_bar = Y.mean()   # 8.7

Sxx = np.sum((X - X_bar) ** 2)           # 82.5
Sxy = np.sum((X - X_bar) * (Y - Y_bar))  # 104.5

beta  = Sxy / Sxx                         # 1.266667
alpha = Y_bar - beta * X_bar              # 1.733333

print(f"X_bar = {X_bar},  Y_bar = {Y_bar}")
print(f"Sxx = {Sxx},  Sxy = {Sxy}")
print(f"beta  (slope)     = {beta:.6f}")
print(f"alpha (intercept) = {alpha:.6f}")
print(f"Fitted line: Y_hat = {alpha:.4f} + {beta:.4f} * X")
assert np.isclose(beta, 19/15)
assert np.isclose(alpha, 26/15)
\end{lstlisting}
\smallskip
\needspace{5\baselineskip}
\noindent\textbf{Out[4]:}
\begin{lstlisting}[style=output]
X_bar = 5.5,  Y_bar = 8.7
Sxx = 82.5,  Sxy = 104.5
beta  (slope)     = 1.266667
alpha (intercept) = 1.733333
Fitted line: Y_hat = 1.7333 + 1.2667 * X
\end{lstlisting}
\vspace{1em}
\begin{tcolorbox}[colback=gray!5!white,colframe=gray!50!black,boxrule=0.5pt,arc=2pt,left=2pt,right=2pt,top=2pt,bottom=2pt]
\textbf{Verification — Method 1: Cramer's Rule}
\end{tcolorbox}
\vspace{1em}
\needspace{5\baselineskip}
\noindent\textbf{In [5]:}
\begin{lstlisting}[style=pycode]
A = np.array([[n, sum_X], [sum_X, sum_X2]], dtype=float)
b = np.array([sum_Y, sum_XY], dtype=float)

det_A       = np.linalg.det(A)
det_alpha   = np.linalg.det(np.column_stack([b,  A[:, 1]]))
det_beta    = np.linalg.det(np.column_stack([A[:, 0], b]))

alpha_cr = det_alpha / det_A
beta_cr  = det_beta  / det_A
print(f"Cramer's Rule → alpha = {alpha_cr:.6f},  beta = {beta_cr:.6f}")
assert np.isclose(alpha_cr, alpha) and np.isclose(beta_cr, beta)
print("[OK] Matches Sxx/Sxy result")
\end{lstlisting}
\smallskip
\needspace{5\baselineskip}
\noindent\textbf{Out[5]:}
\begin{lstlisting}[style=output]
Cramer's Rule → alpha = 1.733333,  beta = 1.266667
[OK] Matches Sxx/Sxy result
\end{lstlisting}
\vspace{1em}
\begin{tcolorbox}[colback=gray!5!white,colframe=gray!50!black,boxrule=0.5pt,arc=2pt,left=2pt,right=2pt,top=2pt,bottom=2pt]
\textbf{Verification — Method 2: Pseudoinverse $(A^T A)^{-1} A^T Y$}
\end{tcolorbox}
\vspace{1em}
\needspace{5\baselineskip}
\noindent\textbf{In [6]:}
\begin{lstlisting}[style=pycode]
A_mat = np.column_stack([np.ones(n), X])   # design matrix [1 | X]
coeffs = np.linalg.inv(A_mat.T @ A_mat) @ A_mat.T @ Y
alpha_ps, beta_ps = coeffs
print(f"Pseudoinverse   → alpha = {alpha_ps:.6f},  beta = {beta_ps:.6f}")
assert np.isclose(alpha_ps, alpha) and np.isclose(beta_ps, beta)
print("[OK] Matches all previous results")
\end{lstlisting}
\smallskip
\needspace{5\baselineskip}
\noindent\textbf{Out[6]:}
\begin{lstlisting}[style=output]
Pseudoinverse   → alpha = 1.733333,  beta = 1.266667
[OK] Matches all previous results
\end{lstlisting}
\vspace{1em}
\begin{tcolorbox}[colback=gray!5!white,colframe=gray!50!black,boxrule=0.5pt,arc=2pt,left=2pt,right=2pt,top=2pt,bottom=2pt]
\textbf{2c. Fitted Line}\par\smallskip
\end{tcolorbox}
\vspace{1em}
\needspace{5\baselineskip}
\noindent\textbf{In [7]:}
\begin{lstlisting}[style=pycode]
Y_hat     = alpha + beta * X
residuals = Y - Y_hat

print(f"Fitted line: Y_hat = {alpha:.4f} + {beta:.4f} * X")
print()
print(f"{'Month':>5} {'X':>5} {'Y':>5} {'Y_hat':>8} {'e_i':>8}")
print("-" * 40)
for i in range(n):
    print(f"  {i+1:2d}   {X[i]:4.1f}  {Y[i]:4.1f}   {Y_hat[i]:7.4f}   {residuals[i]:+7.4f}")
\end{lstlisting}
\smallskip
\needspace{5\baselineskip}
\noindent\textbf{Out[7]:}
\begin{lstlisting}[style=output]
Fitted line: Y_hat = 1.7333 + 1.2667 * X

Month     X     Y    Y_hat      e_i
----------------------------------------
   1    3.0   6.0    5.5333   +0.4667
   2    5.0   9.0    8.0667   +0.9333
   3    2.0   4.0    4.2667   -0.2667
   4    7.0  10.0   10.6000   -0.6000
   5    8.0  12.0   11.8667   +0.1333
   6    1.0   3.0    3.0000   +0.0000
   7    4.0   7.0    6.8000   +0.2000
   8    6.0   8.0    9.3333   -1.3333
   9    9.0  13.0   13.1333   -0.1333
  10   10.0  15.0   14.4000   +0.6000
\end{lstlisting}
\vspace{1em}
\needspace{5\baselineskip}
\noindent\textbf{In [8]:}
\begin{lstlisting}[style=pycode]
# --- Plot 1: Scatter + Regression Line ---
fig1, ax1 = plt.subplots(figsize=(8, 5))
ax1.scatter(X, Y, color='#4F46E5', s=100, edgecolors='white', linewidths=1.5, label='Actual Data')
X_line = np.linspace(0, 13, 100)
ax1.plot(X_line, alpha + beta * X_line, color='#E11D48', linewidth=2.5, label=fr'Regression Line: Y = {alpha:.2f} + {beta:.2f}X')
ax1.scatter([12], [alpha + beta * 12], color='#10B981', s=120, marker='D', edgecolors='white', linewidths=1.5, label='Prediction at X=12')

for i in range(n):
    ax1.annotate(f'({X[i]:.0f}, {Y[i]:.0f})', (X[i], Y[i]), textcoords='offset points', xytext=(8, -12), fontproperties=prop_small, color='#4B5563', bbox=dict(boxstyle="round,pad=0.2", fc="white", ec="none", alpha=0.7))

ax1.set_xlabel('Advertising Budget (X) [$1,000s]', fontproperties=prop_b)
ax1.set_ylabel('Sales (Y) [1,000 units]', fontproperties=prop_b)
ax1.set_title('Advertising Budget vs Sales', fontproperties=prop_title)
legend = ax1.legend(prop=prop_r)
apply_font(ax1)
ax1.set_xlim(-0.5, 13.5)
ax1.set_ylim(0, 18)
plt.tight_layout()
fig1.savefig('plot_1_regression.pdf', dpi=300)
plt.show()
\end{lstlisting}
\noindent\begin{minipage}{\linewidth}
\centering
\vspace{0.4em}
\includegraphics[width=0.96\linewidth]{plot_1_regression.pdf}
\par\vspace{0.6em}
\begin{tcolorbox}[colback=gray!5!white,colframe=gray!50!black,boxrule=0.5pt,arc=2pt,left=6pt,right=6pt,top=5pt,bottom=5pt,unbreakable]
\textbf{Figure 1 (Regression Line)}: Scatter plot displaying the relationship between Advertising Budget and Sales, with the fitted OLS regression line.
\end{tcolorbox}
\vspace{1.0em}
\end{minipage}
\vspace{1em}
\begin{tcolorbox}[colback=gray!5!white,colframe=gray!50!black,boxrule=0.5pt,arc=2pt,left=2pt,right=2pt,top=2pt,bottom=2pt]
\textbf{\large 3. Task 2: Interpretation}\par\smallskip
\textbf{3a. Meaning of Slope and Intercept}\par\smallskip
\begin{itemize}
\item \textbf{Slope ($\beta$ $\approx$ 1.2667):} For each additional \$1,000 in advertising budget, sales are predicted to increase by approximately \textbf{1,267 units}. This is an association in the sample, not proof of causation.
\item \textbf{Intercept ($\alpha$ $\approx$ 1.7333):} At zero advertising budget (X = 0), predicted sales $\approx$ \textbf{1,733 units}. X = 0 is outside the observed range (1–10), so this is an extrapolated baseline.
\end{itemize}
\textbf{3b. Prediction at X = 12 (\$12,000 budget)}\par\smallskip
\end{tcolorbox}
\vspace{1em}
\needspace{5\baselineskip}
\noindent\textbf{In [9]:}
\begin{lstlisting}[style=pycode]
X_pred = 12.0
Y_pred = alpha + beta * X_pred

print(f"X = {X_pred} (i.e., $12,000 advertising budget)")
print(f"Y_hat = {alpha:.4f} + {beta:.4f} * {X_pred} = {Y_pred:.4f} thousand units")
print(f"      ≈ {Y_pred * 1000:,.0f} units")
print()
print("NOTE: X = 12 is beyond the observed maximum (X_max = 10).")
print("This prediction is extrapolation — treat with caution.")
assert np.isclose(Y_pred, 16.9333333333)
\end{lstlisting}
\smallskip
\needspace{5\baselineskip}
\noindent\textbf{Out[9]:}
\begin{lstlisting}[style=output]
X = 12.0 (i.e., $12,000 advertising budget)
Y_hat = 1.7333 + 1.2667 * 12.0 = 16.9333 thousand units
      ≈ 16,933 units

NOTE: X = 12 is beyond the observed maximum (X_max = 10).
This prediction is extrapolation — treat with caution.
\end{lstlisting}
\vspace{1em}
\needspace{5\baselineskip}
\noindent\textbf{In [10]:}
\begin{lstlisting}[style=pycode]
# --- Plot 7: Actual vs Predicted ---
fig7, ax7 = plt.subplots(figsize=(6, 6))
sc = ax7.scatter(Y_hat, Y, c=months, cmap='plasma', s=120, edgecolors='white', linewidths=1.5)
cbar = fig7.colorbar(sc, ax=ax7)
cbar.set_label('Month', fontproperties=prop_b)
cbar.ax.tick_params(labelsize=11)
for t in cbar.ax.get_yticklabels():
    t.set_fontproperties(prop_r)

lims = [min(Y.min(), Y_hat.min()) - 1, max(Y.max(), Y_hat.max()) + 1]
ax7.plot(lims, lims, color='#1F2937', linestyle='--', linewidth=2, label='Perfect Fit')

for i in range(n):
    ax7.annotate(f'M{i+1}', (Y_hat[i], Y[i]), textcoords='offset points', xytext=(-15, 8), fontproperties=prop_small, color='#4B5563', bbox=dict(boxstyle="round,pad=0.2", fc="white", ec="none", alpha=0.6))

ax7.set_xlabel(r'Predicted Sales ($\\hat{Y}$)', fontproperties=prop_b)
ax7.set_ylabel('Actual Sales ($Y$)', fontproperties=prop_b)
ax7.set_title('Actual vs Predicted Sales', fontproperties=prop_title)
ax7.legend(prop=prop_r)
apply_font(ax7)
ax7.set_xlim(lims)
ax7.set_ylim(lims)
ax7.set_aspect('equal')
plt.tight_layout()
fig7.savefig('plot_7_actual_vs_pred.pdf', dpi=300)
plt.show()
\end{lstlisting}
\noindent\begin{minipage}{\linewidth}
\centering
\vspace{0.4em}
\includegraphics[width=0.96\linewidth]{plot_7_actual_vs_pred.pdf}
\par\vspace{0.6em}
\begin{tcolorbox}[colback=gray!5!white,colframe=gray!50!black,boxrule=0.5pt,arc=2pt,left=6pt,right=6pt,top=5pt,bottom=5pt,unbreakable]
\textbf{Figure 7 (Actual vs Predicted)}: Actual vs. Predicted values plot to visualize the overall fit of the model. Points close to the diagonal dashed line indicate accurate predictions.
\end{tcolorbox}
\vspace{1.0em}
\end{minipage}
\vspace{1em}
\begin{tcolorbox}[colback=gray!5!white,colframe=gray!50!black,boxrule=0.5pt,arc=2pt,left=2pt,right=2pt,top=2pt,bottom=2pt]
\textbf{\large 4. Task 3: Model Evaluation}\par\smallskip
\textbf{4a. R² Calculation}\par\smallskip
$$R^2 = 1 - \frac{SS_{\text{res}}}{SS_{\text{tot}}}, \quad SS_{\text{res}} = \sum e_i^2, \quad SS_{\text{tot}} = \sum (Y_i - \bar{Y})^2$$
\end{tcolorbox}
\vspace{1em}
\needspace{5\baselineskip}
\noindent\textbf{In [11]:}
\begin{lstlisting}[style=pycode]
SS_res = np.sum(residuals ** 2)
SS_tot = np.sum((Y - Y_bar) ** 2)
R2     = 1 - SS_res / SS_tot

print(f"SS_res = {SS_res:.6f}")
print(f"SS_tot = {SS_tot:.4f}")
print(f"R²     = {R2:.6f}  ({R2*100:.2f}%)")
print()
print(f"Residual mean = {residuals.mean():.6f}  (≈ 0 → no bias)")
print(f"Residual std  = {residuals.std(ddof=1):.6f}")

assert np.isclose(SS_res, 3.7333333333)
assert np.isclose(SS_tot, 136.1)
assert np.isclose(R2, 0.9725691893)
print("\n[OK] All assertions passed")
\end{lstlisting}
\smallskip
\needspace{5\baselineskip}
\noindent\textbf{Out[11]:}
\begin{lstlisting}[style=output]
SS_res = 3.733333
SS_tot = 136.1000
R²     = 0.972569  (97.26%)

Residual mean = 0.000000  (≈ 0 → no bias)
Residual std  = 0.644061

[OK] All assertions passed
\end{lstlisting}
\vspace{1em}
\begin{tcolorbox}[colback=gray!5!white,colframe=gray!50!black,boxrule=0.5pt,arc=2pt,left=2pt,right=2pt,top=2pt,bottom=2pt]
\textbf{4b. Interpretation}\par\smallskip
97.26\% of the observed variation in sales is explained by the fitted linear relationship with advertising budget.
\textbf{Important caveats:}
\begin{itemize}
\item High R² does \textbf{not} establish that advertising *causes* higher sales.
\item R² alone does not guarantee the model will predict well on new data.
\item Residual mean $\approx$ 0 confirms the model is unbiased over the sample.
\end{itemize}
\end{tcolorbox}
\vspace{1em}
\needspace{5\baselineskip}
\noindent\textbf{In [12]:}
\begin{lstlisting}[style=pycode]
# --- Plot 2: Residuals Lollipop Chart ---
fig2, ax2 = plt.subplots(figsize=(8, 4))
markerline, stemlines, baseline = ax2.stem(X, residuals, linefmt='-', markerfmt='o', basefmt='k-')
plt.setp(markerline, color='#0284C7', markersize=8, markeredgecolor='white', markeredgewidth=1)
plt.setp(stemlines, color='#0284C7', linewidth=2)
plt.setp(baseline, color='#1F2937', linewidth=1.5)

for i in range(n):
    offset = 0.15 if residuals[i] > 0 else -0.35
    ax2.text(X[i], residuals[i] + offset, f'{residuals[i]:.2f}', ha='center', fontproperties=prop_small, color='#0284C7', bbox=dict(boxstyle="round,pad=0.1", fc="white", ec="none", alpha=0.8))

ax2.set_xlabel('Advertising Budget (X) [$1,000s]', fontproperties=prop_b)
ax2.set_ylabel('Residual Error ($e_i$)', fontproperties=prop_b)
ax2.set_title('Residuals (Errors) per Observation', fontproperties=prop_title)
apply_font(ax2)
ax2.set_xlim(-0.5, 12.5)
ax2.set_ylim(min(residuals) - 0.6, max(residuals) + 0.6)
plt.tight_layout()
fig2.savefig('plot_2_residuals_lollipop.pdf', dpi=300)
plt.show()
\end{lstlisting}
\noindent\begin{minipage}{\linewidth}
\centering
\vspace{0.4em}
\includegraphics[width=0.96\linewidth]{plot_2_residuals_lollipop.pdf}
\par\vspace{0.6em}
\begin{tcolorbox}[colback=gray!5!white,colframe=gray!50!black,boxrule=0.5pt,arc=2pt,left=6pt,right=6pt,top=5pt,bottom=5pt,unbreakable]
\textbf{Figure 2 (Residuals Lollipop Plot)}: Lollipop plot illustrating the residuals for each observation, showing the vertical distance between actual and predicted sales.
\end{tcolorbox}
\vspace{1.0em}
\end{minipage}
\vspace{1em}
\needspace{5\baselineskip}
\noindent\textbf{In [13]:}
\begin{lstlisting}[style=pycode]
# --- Plot 3: Residuals vs Fitted ---
fig3, ax3 = plt.subplots(figsize=(6, 4.5))
ax3.scatter(Y_hat, residuals, color='#F59E0B', s=90, edgecolors='white', linewidths=1.2)
ax3.axhline(0, color='#1F2937', linestyle='--', linewidth=1.5)
for i in range(n):
    ax3.annotate(f'M{i+1}', (Y_hat[i], residuals[i]), textcoords='offset points', xytext=(6, 5), fontproperties=prop_small, color='#4B5563')
ax3.set_xlabel(r'Fitted Values ($\\hat{Y}$)', fontproperties=prop_b)
ax3.set_ylabel('Residuals', fontproperties=prop_b)
ax3.set_title('Residuals vs Fitted Values', fontproperties=prop_title)
apply_font(ax3)
plt.tight_layout()
fig3.savefig('plot_3_res_vs_fitted.pdf', dpi=300)
plt.show()
\end{lstlisting}
\noindent\begin{minipage}{\linewidth}
\centering
\vspace{0.4em}
\includegraphics[width=0.96\linewidth]{plot_3_res_vs_fitted.pdf}
\par\vspace{0.6em}
\begin{tcolorbox}[colback=gray!5!white,colframe=gray!50!black,boxrule=0.5pt,arc=2pt,left=6pt,right=6pt,top=5pt,bottom=5pt,unbreakable]
\textbf{Figure 3 (Residuals vs Fitted)}: Residuals vs. Fitted plot to check for non-linear patterns and homoscedasticity. Points scattered randomly around the 0 line indicate a good fit.
\end{tcolorbox}
\vspace{1.0em}
\end{minipage}
\vspace{1em}
\needspace{5\baselineskip}
\noindent\textbf{In [14]:}
\begin{lstlisting}[style=pycode]
# --- Plot 4: Residual Distribution (Histogram) ---
fig4, ax4 = plt.subplots(figsize=(6, 4.5))
ax4.hist(residuals, bins=6, color='#8B5CF6', edgecolor='white', alpha=0.8, density=True)
x_range = np.linspace(residuals.min() - 0.5, residuals.max() + 0.5, 100)
ax4.plot(x_range, stats.norm.pdf(x_range, 0, residuals.std(ddof=1)), color='#E11D48', linewidth=2)
ax4.axvline(0, color='#1F2937', linestyle='--', linewidth=1.2)
ax4.set_xlabel('Residual Value', fontproperties=prop_b)
ax4.set_ylabel('Density', fontproperties=prop_b)
ax4.set_title('Residuals Distribution', fontproperties=prop_title)
apply_font(ax4)
plt.tight_layout()
fig4.savefig('plot_4_res_dist.pdf', dpi=300)
plt.show()
\end{lstlisting}
\noindent\begin{minipage}{\linewidth}
\centering
\vspace{0.4em}
\includegraphics[width=0.96\linewidth]{plot_4_res_dist.pdf}
\par\vspace{0.6em}
\begin{tcolorbox}[colback=gray!5!white,colframe=gray!50!black,boxrule=0.5pt,arc=2pt,left=6pt,right=6pt,top=5pt,bottom=5pt,unbreakable]
\textbf{Figure 4 (Residuals Distribution)}: Histogram of residuals to assess whether they approximate a normal distribution (Normality assumption).
\end{tcolorbox}
\vspace{1.0em}
\end{minipage}
\vspace{1em}
\needspace{5\baselineskip}
\noindent\textbf{In [15]:}
\begin{lstlisting}[style=pycode]
# --- Plot 5: Normal Q-Q Plot ---
fig5, ax5 = plt.subplots(figsize=(6, 4.5))
(osm, osr), (slope_qq, intercept_qq, _) = stats.probplot(residuals, dist='norm')
ax5.scatter(osm, osr, color='#EC4899', s=90, edgecolors='white', linewidths=1.2)
xq = np.array([osm.min(), osm.max()])
ax5.plot(xq, slope_qq * xq + intercept_qq, color='#1F2937', linewidth=2)
ax5.set_xlabel('Theoretical Quantiles', fontproperties=prop_b)
ax5.set_ylabel('Sample Quantiles', fontproperties=prop_b)
ax5.set_title('Normal Q-Q Plot', fontproperties=prop_title)
apply_font(ax5)
plt.tight_layout()
fig5.savefig('plot_5_qq.pdf', dpi=300)
plt.show()
\end{lstlisting}
\noindent\begin{minipage}{\linewidth}
\centering
\vspace{0.4em}
\includegraphics[width=0.96\linewidth]{plot_5_qq.pdf}
\par\vspace{0.6em}
\begin{tcolorbox}[colback=gray!5!white,colframe=gray!50!black,boxrule=0.5pt,arc=2pt,left=6pt,right=6pt,top=5pt,bottom=5pt,unbreakable]
\textbf{Figure 5 (Q-Q Plot)}: Q-Q plot comparing the quantiles of the residuals against a standard normal distribution. Points closely following the diagonal line suggest normality.
\end{tcolorbox}
\vspace{1.0em}
\end{minipage}
\vspace{1em}
\needspace{5\baselineskip}
\noindent\textbf{In [16]:}
\begin{lstlisting}[style=pycode]
# --- Plot 6: Scale-Location Plot ---
fig6, ax6 = plt.subplots(figsize=(6, 4.5))
sqrt_abs_res = np.sqrt(np.abs(residuals))
ax6.scatter(Y_hat, sqrt_abs_res, color='#10B981', s=90, edgecolors='white', linewidths=1.2)
for i in range(n):
    ax6.annotate(f'M{i+1}', (Y_hat[i], sqrt_abs_res[i]), textcoords='offset points', xytext=(6, 5), fontproperties=prop_small, color='#4B5563')
ax6.set_xlabel(r'Fitted Values ($\\hat{Y}$)', fontproperties=prop_b)
ax6.set_ylabel(r'$\\sqrt{|Residuals|}$', fontproperties=prop_b)
ax6.set_title('Scale-Location Plot', fontproperties=prop_title)
apply_font(ax6)
plt.tight_layout()
fig6.savefig('plot_6_scale_loc.pdf', dpi=300)
plt.show()
\end{lstlisting}
\noindent\begin{minipage}{\linewidth}
\centering
\vspace{0.4em}
\includegraphics[width=0.96\linewidth]{plot_6_scale_loc.pdf}
\par\vspace{0.6em}
\begin{tcolorbox}[colback=gray!5!white,colframe=gray!50!black,boxrule=0.5pt,arc=2pt,left=6pt,right=6pt,top=5pt,bottom=5pt,unbreakable]
\textbf{Figure 6 (Scale-Location)}: Scale-Location plot to check for homoscedasticity (equal variance of residuals) across the range of fitted values. The trend line should be relatively horizontal.
\end{tcolorbox}
\vspace{1.0em}
\end{minipage}

\end{document}
